\chapter{Identités et fonctions trigonométriques}
\section*{Exercice 15.1 -- Addition d'angles}
Calculer exactement $\sin(\pi/12)$ à partir de $\pi/4-\pi/6$.

\section*{Exercice 15.2 -- Angle double}
Démontrer $1-\cos(2x)=2\sin^2x$ puis résoudre $\cos(2x)=1$.

\section*{Exercice 15.3 -- Dériver un signal}
Pour $f(t)=3\cos(2t+\pi/6)$, donner amplitude, période, dérivée et dérivée seconde.

\section*{Exercice 15.4 -- Rotation}
Appliquer la rotation d'angle $\pi/2$ au vecteur $(2,-3)$ et vérifier sa norme.

\section*{Exercice 15.5 -- Produit en somme}
Prouver $2\sin a\sin b=\cos(a-b)-\cos(a+b)$ et l'appliquer à $\sin^2x$.

\section*{Exercice 15.6 -- Limite fondamentale}
À partir de $\cos h\le\sin h/h\le1$ pour $0<h<\pi/2$, justifier la limite en $0$ des deux côtés.
